\documentclass[12pt]{article}
\usepackage[utf8]{inputenc}
\usepackage[T1]{fontenc}
\usepackage[margin=1in]{geometry}
\usepackage{setspace}
\usepackage{xcolor}
\usepackage[colorlinks=true, urlcolor=blue!60!black, linkcolor=blue!60!black]{hyperref}
\setlength{\parindent}{0em}
\setlength{\parskip}{0.45em}
\renewcommand{\baselinestretch}{1.0}
\begin{document}
\begin{center}
{\Large \textit{Board of Trustees of the University of Alabama v. Garrett}, 531 U.S. 356 (2001)\par}
\vspace{0.3em}
{\normalsize Case Brief --- One-Page Edition\par}
\end{center}

\textbf{Memory Jogger.} No ADA Title I damages suits against States: Congress showed no pattern of unconstitutional state discrimination and its remedy was not congruent and proportional (\href{https://www.courtlistener.com/opinion/118407/board-of-trustees-of-univ-of-ala-v-garrett/}{CourtListener}).

\textbf{Facts.} Nurse Patricia Garrett lost her Director of Nursing post at a state university hospital after breast-cancer treatment; security officer Milton Ash was denied accommodations for asthma and sleep apnea. Both sued Alabama for damages under Title I, which expressly abrogates state immunity.

\textbf{Procedural History.} The district court granted summary judgment for the State; the Eleventh Circuit reversed.

\textbf{Votes.} Reversed 5--4. Rehnquist wrote for the Court; Kennedy (with O'Connor) concurred; Breyer (with Stevens, Souter, and Ginsburg) dissented.

\textbf{Holding.} The Eleventh Amendment bars private damages suits against States under Title I. For rights subject to rational-basis review, Congress must document a pattern of unconstitutional discrimination by the States themselves and enact congruent and proportional remedies.

\textbf{Reasoning.} Under \textit{Cleburne}, disability classifications need only be rational. Excluding local governments, the record held only ``half a dozen examples'' of state misconduct, and Title I's accommodation and disparate-impact rules far exceed equal protection. Kennedy demanded ``confirming judicial documentation'' before accusing sovereign States. Breyer replied that ``a legislature is not a court of law'' and that rational-basis deference restrains courts, not Congress. Title I still binds States through federal suits and \textit{Ex parte Young} injunctions.

\textbf{Examples.} \emph{Same Side:} genetic-status damages suits on a private-sector record are barred. \emph{Opposite Side:} courthouse-access suits survive (\textit{Tennessee v. Lane}). \emph{Fence-Sitter:} suits over inaccessible state-university websites.

\textbf{Critique.} Critics see judicial activism that makes Section 5 protection of non-suspect classes nearly impossible, and textualists note the Eleventh Amendment does not reach suits by a State's own citizens.

\textbf{Key Quotation.} ``States are not required by the Fourteenth Amendment to make special accommodations for the disabled, so long as their actions toward such individuals are rational. They could quite hardheadedly---and perhaps hardheartedly---hold to job-qualification requirements which do not make allowance for the disabled.'' --- Chief Justice Rehnquist, 531 U.S. at 367.

\end{document}
